% Imporditud: tools/import_eksperimendid.py
% Allikas: Eesti füüsikaolümpiaadi eksperimendi ülesannete
%   kogu (Rael Kalda, 2023), ülesanne Ü171, lahendus L171.
\setAuthor{Jaan Kalda}
\setRound{lõppvoor}
\setYear{2015}
\setNumber{G E2}
\setDifficulty{5}
\setTopic{elektriahelad}
\setVahendid{kaks voltmeetrit, kaks takistit ($R_1$ = 1 k$\Omega$ ja $R_2$ = 10 k$\Omega$), kondensaator (mahtuvus suurusjärgus 20 mF), patarei pingega 1,5 V, juhtmed, millimeeterpaber}
\setPunktid{14}
\setAllikas{Eksperimendi kogu Ü171, lahendus lk 130}
\setLahendusseis{toores}

\prob{Valgusdiood}
Määrata valgusdioodi voolutugevuse sõltuvus pingest vahemikus 0 V$^{-3}$ V (joonistada graafik).

\hint

\solu
Ühendame patarei järjestikku takistiga takistusega R ja dioodiga ning mõõdame pinget takistil UT ja pinget dioodil UD. Voolutugevus läbi dioodi on I = UT /R. Näeme, et patarei pingest ei piisa selleks, et saada läbi dioodi nõutavat voolutugevust. Seepärast laeme kondensaatori ja ühendame selle järjestikku patareiga, nii et pinge kahekordistub. Alustame takistiga takistusega R = 1 k$\Omega$, sest R = 10 k$\Omega$ puhul ei saavutaks me soovitud suurimat voolutugevust (100 $\mu$A). Vajalikud andmed saame tühjenemise käigus pingeid samaaegselt registreerides. Umbes I = 10 $\mu$A juures (täpne hetk pole tähtis) vahetame vooluringis takistit, sest muidu lähevad takistil mõõdetavad pinged nii väikeseks, et voolu mõõtmiseks kasutatava voltmeetri näidatavate tüvenumbrite arv jääb liiga napiks. Väga väikeste voolude juures (alla I = 5 $\mu$A) tühjeneb kondensaator nii aeglaselt, et voolu vähenemist ei jõua ära oodata. Siis tuleb teise takisti abil kondensaatorit järg-järgult tühjendada (st ühendada see mahtuvusega rööbiti). Sobilikku tühjenemisastet saab samal ajal jälgida dioodivoolu abil (st vaadata pinget dioodiga jadamisi oleval takistil)
\probend
